\documentclass[aspectratio=169,10pt]{beamer}
\usepackage[T1]{fontenc}
\usepackage[utf8]{inputenc}
\usepackage[spanish,es-noshorthands]{babel}
\usepackage{lmodern}
\usepackage{tikz}
\usepackage{array,booktabs,tabularx}
\usetikzlibrary{positioning,arrows.meta}
\definecolor{ink}{HTML}{142F3D}
\definecolor{teal}{HTML}{087C78}
\definecolor{mint}{HTML}{EAF7F2}
\definecolor{gold}{HTML}{E49A3F}
\definecolor{light}{HTML}{F4F7F8}
\definecolor{muted}{HTML}{536976}
\setbeamercolor{normal text}{fg=ink,bg=white}
\setbeamercolor{frametitle}{fg=ink,bg=white}
\setbeamerfont{frametitle}{size=\Large,series=\bfseries}
\setbeamertemplate{navigation symbols}{}
\setbeamertemplate{headline}{}
\setbeamertemplate{footline}{\leavevmode\hbox{\begin{beamercolorbox}[wd=.88\paperwidth,ht=3.4ex,dp=1ex,leftskip=.6cm]{normal text}\scriptsize\color{teal}\textbf{FARO\textit{CX}}\quad\color{muted} | Intelligence for retention \qquad\color{muted} Kevin Andrés Erazo Yepez\end{beamercolorbox}\begin{beamercolorbox}[wd=.12\paperwidth,ht=3.4ex,dp=1ex,center]{normal text}\scriptsize\color{teal}\insertframenumber\,/\,5\end{beamercolorbox}}}
\newcommand{\eyebrow}[1]{{\vspace{-.15cm}\scriptsize\bfseries\color{teal}\MakeUppercase{#1}\par\vspace{.18cm}}}
\newcommand{\kpi}[3]{\begin{minipage}{.29\linewidth}\colorbox{mint}{\parbox{\dimexpr\linewidth-2\fboxsep}{\centering\vspace{.16cm}\LARGE\bfseries\color{teal}#1\par\vspace{.09cm}{\small\color{ink}#2}\vspace{.16cm}}}\par\vspace{.08cm}{\scriptsize\color{muted}#3}\end{minipage}}
\newcommand{\card}[2]{\colorbox{light}{\parbox{.90\linewidth}{\vspace{.12cm}\textbf{\color{teal}#1}\par\vspace{.09cm}{\small\color{ink}#2}\vspace{.14cm}}}}
\begin{document}

% 1. PROBLEMA Y POSICIONAMIENTO
\begin{frame}[t]{La fuga se anticipa. La decisión también.}
\eyebrow{01 / Oportunidad comercial}
\vspace{.07cm}
{\large En una operadora, cada contacto cuesta. Actuar a ciegas desperdicia presupuesto; no actuar deja pasar oportunidades de retención.}\par
\vspace{.28cm}
\kpi{7.043}{clientes}{Cartera histórica Telco}\hfill
\kpi{26,5\%}{abandono observado}{No es una predicción futura}\hfill
\kpi{1.869}{bajas registradas}{Base pública de demostración}
\vspace{.40cm}
\colorbox{ink}{\parbox{.96\linewidth}{\vspace{.19cm}\centering\large\bfseries\color{white}No preguntamos solo quién podría irse.\\[.05cm]Preguntamos: ¿en quién conviene invertir primero?\vspace{.19cm}}}
\vspace{.21cm}
{\scriptsize\color{muted}FaroCX convierte un pronóstico estadístico en un escenario de priorización comercial; no garantiza que una intervención evite una baja.}
\end{frame}

% 2. LA SOLUCION
\begin{frame}[t]{FaroCX: riesgo, explicación y acción en un solo flujo}
\eyebrow{02 / Propuesta diferencial}
\vspace{.10cm}
\begin{columns}[T,onlytextwidth]
\column{.31\textwidth}
\card{01 \quad DETECTAR}{Estima el riesgo individual con un modelo evaluado sobre datos independientes.}
\vspace{.20cm}
\card{02 \quad ENTENDER}{SHAP y LIME muestran qué señales influyen en cada pronóstico.}
\column{.31\textwidth}
\card{03 \quad PRIORIZAR}{Ordena candidatos según el beneficio neto esperado y el presupuesto disponible.}
\vspace{.20cm}
\card{04 \quad MEDIR}{Propone validar el impacto de las campañas con un piloto controlado A/B.}
\column{.32\textwidth}
{\color{teal}\bfseries El valor no es una lista de porcentajes.}\par
\vspace{.18cm}
{\large Es una \textbf{lista de decisiones justificables}.}\par
\vspace{.22cm}
{\small\color{muted}Riesgo $\times$ eficacia supuesta $\times$ margen estimado $-$ costo de contacto.}
\end{columns}
\vspace{.32cm}
{\small\color{teal}\textbf{Diferenciación:} el equipo de CRM puede variar presupuesto y costos para comparar escenarios antes de contactar.}
\end{frame}

% 3. DEMO
\begin{frame}[t]{Una demostración que termina en una decisión}
\eyebrow{03 / Producto funcionando en la nube}
\vspace{.10cm}
{\large Tres pantallas; una pregunta de negocio: \textbf{¿a quién contactamos primero?}}\par
\vspace{.28cm}
\begin{columns}[T,onlytextwidth]
\column{.32\textwidth}
\card{PERFIL}{Introducir datos de un cliente y consultar el riesgo estimado de abandono.}
\column{.32\textwidth}
\card{EXPLICACIÓN}{Mostrar factores locales SHAP y LIME: por qué el modelo emitió esa alerta.}
\column{.32\textwidth}
\card{CAMPAÑA}{Ajustar presupuesto, margen y costo de contacto; exportar una cartera priorizada.}
\end{columns}
\vspace{.40cm}
\colorbox{mint}{\parbox{.96\linewidth}{\vspace{.16cm}\centering\bfseries\color{ink}FaroCX hace visible el paso de \textcolor{teal}{dato} a \textcolor{teal}{decisión}.\vspace{.16cm}}}
\vspace{.24cm}
{\scriptsize\color{muted}Los escenarios económicos son hipotéticos. Esta versión usa datos públicos de demostración y no integra información privada de clientes reales.}
\end{frame}

% 4. EVIDENCIA
\begin{frame}[t]{Evidencia comprobable, sin promesas infladas}
\eyebrow{04 / Validación e interpretabilidad}
\vspace{.03cm}
\begin{columns}[T,onlytextwidth]
\column{.50\textwidth}
{\small\bfseries Rendimiento en test reservado}\par\vspace{.08cm}
\begin{tabular}{@{}lcc@{}}\toprule
\small Modelo & \small F1 & \small ROC-AUC\\\midrule
\small Logistic Regression & \textbf{0.604} & \textbf{0.842}\\
\small RF seleccionado por NSGA-II & 0.581 & 0.830\\\bottomrule
\end{tabular}
\vspace{.17cm}
{\scriptsize\color{muted}Logistic Regression obtuvo mejores métricas en test. El Random Forest fue seleccionado antes mediante criterios de F1-CV y latencia (F1-CV=0.5608).}
\column{.46\textwidth}
{\small\bfseries Modelos explicables para CRM}\par\vspace{.09cm}
{\small SHAP global destaca: \textbf{Contract, tenure, OnlineSecurity}.}\par\vspace{.15cm}
{\small LIME y SHAP explican el perfil individual.}\par\vspace{.15cm}
{\small Grid Search, Random Search, Hyperband y NSGA-II documentan el proceso de optimización.}
\end{columns}
\vspace{.32cm}
\colorbox{light}{\parbox{.96\linewidth}{\vspace{.12cm}\scriptsize\color{muted}Random Search y Hyperband: F1-CV 0.5736; Random Search fue más rápido en esta ejecución.\par\vspace{.06cm}Ni las explicaciones ni las simulaciones establecen causalidad ni efectividad comercial.\vspace{.12cm}}}
\end{frame}

% 5. NEGOCIACION
\begin{frame}[t]{No le pido que compre una promesa. Le propongo medirla.}
\eyebrow{05 / Propuesta comercial}
\vspace{.08cm}
\begin{columns}[T,onlytextwidth]
\column{.48\textwidth}
{\large\bfseries Piloto acotado con su equipo de CRM}\par\vspace{.18cm}
\card{ALCANCE}{Cartera autorizada, perfiles de riesgo, explicaciones y escenarios de contacto.}
\vspace{.16cm}
\card{MODELO COMERCIAL}{Referencia inicial: USD 59/mes (hipótesis ajustable al alcance).}
\column{.48\textwidth}
{\large\bfseries Criterios de continuidad}\par\vspace{.18cm}
\begin{enumerate}\small
\item Calidad predictiva sobre datos del cliente.
\item Costo por contacto y adopción por CRM.
\item Bajas incrementales evitadas, medidas con control A/B.
\item Privacidad, sesgos y seguridad revisados.
\end{enumerate}
\end{columns}
\vspace{.33cm}
\colorbox{ink}{\parbox{.96\linewidth}{\vspace{.21cm}\centering\large\bfseries\color{white}¿Definimos un piloto con presupuesto, métricas y fecha de revisión?\vspace{.21cm}}}
\vspace{.13cm}
{\scriptsize\color{muted}Se necesita un horizonte adecuado para observar bajas; no se afirma ROI ni reducción de churn demostrados.}
\end{frame}
\end{document}
